\Qhead{3.2}{Корреляционные и спектральные характеристики стационарных случайных процессов}{Корреляционные и спектральные характеристики стационарных случайных процессов. Теорема Винера-Хинчина. Белый шум и другие примеры спектров и корреляционных функций.}
\Subfile{3.2.1}{sections/03_statistics/q02/s01}

\recap{Корреляция и спектр — взаимные преобразования Фурье. Спектр неотрицателен и интегрируется в дисперсию; время корреляции обратно связано с шириной спектра.}
\board{Пара преобразований, конечновременной вывод, белый и лоренцев шум; переход от $S(\omega)$ к $G(f)$.}
\oral{Почему нельзя приписать идеальному белому шуму конечную мгновенную дисперсию?}
\bait{Можно оставить переход к одностороннему спектру: множитель два возникает из объединения положительных и отрицательных частот.}
\sources{Биленко и соавт.~\cite{bilenko}, §§4.1.4--4.1.7, PDF-с.~150--155. Соглашение о преобразовании Фурье приведено к единой нормировке сборника.}
